% =====================================================================
% Guía de exposición del Plan de Tesis (PFC II), 7 minutos
% Acompaña a Diapositivas.tex. Explica solo lo que se ve en cada
% diapositiva, en lenguaje simple y con la forma de decir cada símbolo.
% Compilar con: latexmk -pdf Guia_Exposicion.tex
% =====================================================================
\documentclass[11pt,a4paper]{article}
\usepackage[utf8]{inputenc}
\usepackage[T1]{fontenc}
\usepackage{lmodern}
\usepackage[spanish,es-noshorthands]{babel}
\usepackage{amsmath,amssymb}
\usepackage{booktabs,array,longtable}
\usepackage[margin=2.2cm]{geometry}
\usepackage{xcolor}
\usepackage[colorlinks=true,linkcolor=blue!50!black]{hyperref}

\definecolor{ucspblue}{HTML}{14325A}
\definecolor{acc}{HTML}{C8651B}

\setlength{\parindent}{0pt}
\setlength{\parskip}{0.6ex}
\renewcommand{\arraystretch}{1.2}

% Cuadros "qué decir" (azul) y "si te preguntan" (naranja)
\setlength{\fboxsep}{6pt}
\setlength{\fboxrule}{0.6pt}
\newsavebox{\cajabox}
\newcommand{\cajainicio}[3]{\def\cajafc{#1}\def\cajabc{#2}\par\medskip\noindent
  \begin{lrbox}{\cajabox}\begin{minipage}{\dimexpr\linewidth-2\fboxsep-2\fboxrule\relax}%
  \small\textbf{\color{#1}#3}\par\smallskip}
\newcommand{\cajafin}{\end{minipage}\end{lrbox}\fcolorbox{\cajafc}{\cajabc}{\usebox{\cajabox}}\par\medskip}
\newenvironment{guion}{\cajainicio{ucspblue}{ucspblue!6}{Qué decir}}{\cajafin}
\newenvironment{preg}[1]{\cajainicio{acc}{acc!7}{Si te preguntan: #1}}{\cajafin}

% Tabla de elementos de la diapositiva
\newenvironment{partes}{%
  \begin{longtable}{@{}>{\raggedright\arraybackslash}p{4.3cm}>{\raggedright\arraybackslash}p{6.0cm}>{\raggedright\arraybackslash}p{5.0cm}@{}}
  \toprule \textbf{Lo que se ve (y cómo se dice)} & \textbf{Qué es, en simple} & \textbf{Ejemplo}\\ \midrule \endhead}{%
  \bottomrule\end{longtable}}

% Cómo se dice un símbolo o sigla
\newcommand{\se}[1]{\newline{\small\color{ucspblue}(se dice «#1»)}}

\newcommand{\BBAC}{BBAC$_I$}
\newcommand{\STCC}{ST-CC}
\newcommand{\ASTCC}{AST-CC}

\title{\textbf{Guía de exposición}\\[0.5ex]\large Extensión de ST-CC con refinamiento jerárquico por bloques\\para la detección de patrones espacio-temporales locales}
\author{Luis Fabian Diaz Alvizuri}
\date{Arequipa, octubre de 2026}

\begin{document}
\maketitle

\section*{Cómo usar esta guía}

Cada sección corresponde a una diapositiva, en el orden en que se expone. Tiene un cuadro azul con lo que conviene decir, una tabla que explica en lenguaje simple cada cosa que se ve en la diapositiva (con la forma de decir cada símbolo entre paréntesis) y cuadros naranjas con respuestas cortas a preguntas probables. No hay demostraciones ni deducciones, solo definiciones y ejemplos.

\begin{center}
\small
\begin{tabular}{@{}clcl@{}}
\toprule
N.º & Diapositiva & Tiempo & Mensaje que debe quedar \\
\midrule
1 & Portada & 10 s & Quién soy y de qué trata. \\
2 & Contexto: ST-CC & 35 s & ST-CC da un tablero global. \\
3 & La limitación & 35 s & Un foco local obliga a refinar todo. \\
4 & Planteamiento & 40 s & Tres preguntas sin respuesta estadística. \\
5 & Objetivos & 30 s & Un general y cuatro específicos. \\
6 & Estado del arte & 35 s & Tres familias, ninguna refina por bloque. \\
7 & Brecha & 25 s & Nadie reúne las tres características. \\
8 & AST-CC & 70 s & Refinar mientras haya algo real y relevante. \\
9 & Regla de decisión & 50 s & Qué se pregunta a cada bloque. \\
10 & Calibración y parada & 55 s & Cómo se evita el falso refinamiento. \\
11 & Resultados esperados & 35 s & Cómo se va a evaluar. \\
\midrule
 & \textbf{Total} & \textbf{7 min} & \\
\bottomrule
\end{tabular}
\end{center}

\section*{Cómo se leen los símbolos}

Estos símbolos aparecen en varias diapositivas. Si se te olvida uno durante la exposición, esta tabla te dice cómo pronunciarlo.

\begin{center}
\small
\begin{longtable}{@{}>{\centering\arraybackslash}p{2.6cm}>{\raggedright\arraybackslash}p{4.8cm}>{\raggedright\arraybackslash}p{8.0cm}@{}}
\toprule
\textbf{Símbolo} & \textbf{Se dice} & \textbf{Para qué sirve en la tesis} \\
\midrule
\endhead
$\alpha$ & «alfa» & Nivel de significancia: cuánto error se tolera. \\
$\delta$ & «delta» & Umbral mínimo de tamaño de efecto: qué tan grande debe ser un patrón para que valga la pena. \\
$\mu$ & «mu» & La media (el promedio esperado) de una celda. \\
$\hat\varphi$ & «fi sombrero» & Factor que mide cuánto más variable es el dato que lo que permite Poisson. \\
$\sum$ & «sumatoria» o «suma de» & Sumar una cantidad sobre todos los elementos indicados. \\
$\tilde{\ }$ (en $\tilde I_B$) & «tilde» (el signo ``virgulilla'') & Indica que la cantidad ya está corregida por el ruido. \\
$G^2_B$ & «G cuadrado de B» & Medida de cuánto patrón interno tiene el bloque $B$. \\
$G^2_{B,s}$ & «G cuadrado de B, s» & Parte de ese patrón que explica el corte $s$. \\
$\tilde I_B$ & «I tilde de B» & Tamaño de efecto del bloque $B$. \\
$\tilde g_{B,s}$ & «g tilde de B, s» & Ganancia del corte $s$ en el bloque $B$. \\
$\alpha_B$ & «alfa sub B» & Nivel de significancia que le toca al bloque $B$. \\
$\alpha_C=\alpha_B\,N_C/N_B$ & «alfa sub C es igual a alfa sub B por N sub C sobre N sub B» & El nivel de un hijo es el del padre, repartido según los viajes de cada hijo. \\
$K_R\times K_T$ & «K sub R por K sub T» & Tamaño del tablero: grupos de filas por grupos de columnas. \\
$X_1,\dots,X_D$ & «X sub uno, hasta X sub D» & Las matrices de cada uno de los $D$ días. \\
$\ge$ & «mayor o igual que» & Comparación (por ejemplo $\tilde I_B\ge\delta$). \\
$2\times2$ & «dos por dos» & Un corte que parte filas y columnas a la vez. \\
$\checkmark$ y $\times$ & «cumple» y «no cumple» & Marcas de la tabla de brecha. \\
$a\,\mu^b$ & «a por mu a la b» & Fórmula de cuánto crece la variabilidad con el volumen. \\
\midrule
ST-CC & «ese te, ce ce» & \emph{Spatio-Temporal Co-Clustering}, método de Liu et al. (2021). \\
AST-CC & «a, ese te, ce ce» & El método propuesto en esta tesis. \\
BBAC$_I$ & «be be a ce, sub i» & Algoritmo de co-agrupamiento que usa ST-CC. \\
FWER & «efe, doble u, e, erre» & Probabilidad de al menos un falso refinamiento en todo el árbol. \\
ARI & «a, erre, i» & Medida de qué tan parecidas son dos particiones. \\
BIC & «bic» & Medida de ajuste que penaliza los parámetros. \\
OE1 a OE4 & «o e uno» a «o e cuatro» & Objetivos específicos. \\
\bottomrule
\end{longtable}
\end{center}

\subsection*{Un bloque de ejemplo que se reutiliza}

Un bloque con 4 regiones (filas) y 2 horarios (columnas), con viajes así:

\begin{center}
\begin{tabular}{@{}lcc@{}}
\toprule
 & Mañana & Noche \\
\midrule
Región 1 & 30 & 10 \\
Región 2 & 28 & 12 \\
Región 3 & 10 & 30 \\
Región 4 & 12 & 28 \\
\bottomrule
\end{tabular}
\end{center}

Las regiones 1 y 2 son «de mañana» y las 3 y 4 son «de noche». Si todas se comportaran igual, cada celda tendría 20 viajes. Un tablero como el de ST-CC resumiría este bloque con un solo promedio y perdería ese patrón. Su valor de $G^2_B$ es unos 34, que es alto y significa que hay patrón. Un corte de filas que separe $\{1,2\}$ de $\{3,4\}$ explica casi todo ese patrón (unos 33.6 de 34.1) y deja dos hijos casi uniformes.

% =====================================================================
\section{Portada}

\begin{guion}
Buenas tardes. Soy Luis Fabian Diaz Alvizuri y presento mi plan de tesis, asesorado por el Dr. Erick Gómez Nieto. La tesis propone una extensión de ST-CC que refina localmente sus bloques para detectar patrones espacio-temporales locales.
\end{guion}

\begin{preg}{¿Qué significa ST-CC?}
Son las siglas de \emph{Spatio-Temporal Co-Clustering}. Es el método de Liu et al. (2021) que agrupa a la vez las regiones y los intervalos de tiempo de una matriz de viajes.
\end{preg}

% =====================================================================
\section{Contexto: ST-CC y el co-agrupamiento en tablero}

\begin{guion}
Datos como viajes en taxi describen dónde y cuándo ocurre un fenómeno. Liu et al. proponen ST-CC, que agrupa filas y columnas de una matriz región por intervalo con el algoritmo BBAC$_I$. El resultado es un tablero de $K_R\times K_T$ bloques. En Manhattan usaron unos 10 millones de viajes, 66 regiones y 12 intervalos de 2 horas. Es fácil de interpretar, pero su estructura es global.
\end{guion}

\begin{partes}
Dónde y cuándo & El fenómeno tiene una parte espacial (la región) y una temporal (el horario), y ambas están relacionadas. & Hay más viajes en Midtown a las 8 de la mañana que a las 3 de la madrugada. \\
Co-agrupamiento (\emph{co-clustering}) \se{co-agrupamiento} & Agrupar a la vez las filas y las columnas de una matriz, no solo una de las dos. & Se juntan regiones parecidas y, a la vez, horarios parecidos. \\
ST-CC \se{ese te, ce ce} & Método de Liu et al. (2021) que hace co-agrupamiento sobre una matriz de regiones por tiempo. & Es el método del que parte la tesis. \\
Matriz regiones $\times$ intervalos de tiempo \se{matriz de regiones por intervalos de tiempo} & Tabla donde cada fila es una región, cada columna es un intervalo y cada celda cuenta viajes. & Celda (región A, 8 a 10 h) $=$ 120 viajes. \\
BBAC$_I$ \se{be be a ce, sub i} & Algoritmo que reparte las filas y las columnas en grupos para que cada bloque resultante sea lo más uniforme posible. La $I$ indica la medida de error que usa, la I-divergencia. & Es el algoritmo que produce el tablero de ST-CC. \\
Tablero de $K_R\times K_T$ bloques \se{K sub R por K sub T} & $K_R$ es el número de grupos de filas y $K_T$ el de grupos de columnas. Juntos forman una cuadrícula de bloques rectangulares. & Con $K_R=3$ y $K_T=3$ hay 9 bloques. \\
Unos 10 millones de viajes, 66 regiones, 12 intervalos de 2 horas & Los datos de Manhattan en noviembre de 2017 que usaron Liu et al. & $66\times12=792$ celdas. \\
Estructura global & Los cortes entre grupos cruzan toda la matriz de lado a lado. & Una línea entre dos grupos de regiones atraviesa todos los horarios. \\
Cada bloque se resume con un valor (figura) & En el tablero, cada bloque se pinta de un solo color: su promedio. & Un bloque con celdas de 10 y de 70 se pinta como 30. \\
\end{partes}

\begin{preg}{¿Qué es la I-divergencia?}
Es la medida de error que usa BBAC$_I$ para decidir qué tan bien un tablero representa los datos. Con conteos equivale a una medida estándar de ajuste de la estadística (la devianza de Poisson), y por eso permite usar pruebas de hipótesis.
\end{preg}

\begin{preg}{¿Qué hace ST-CC después de formar el tablero?}
Reagrupa los bloques con K-means según su promedio y su dispersión, en 4 niveles de actividad. Eso fusiona bloques parecidos, pero no divide ninguno.
\end{preg}

\begin{preg}{¿Por qué conteos y no promedios por día?}
Porque con conteos la medida de error de BBAC$_I$ coincide con la de Poisson, y eso es lo que permite calibrar las pruebas.
\end{preg}

% =====================================================================
\section{La limitación: un detalle local obliga a refinar todo}

\begin{guion}
En el panel (a) hay un foco local dentro de un bloque. En el panel (b), el tablero de ST-CC lo promedia y desaparece. Como cada corte de filas o columnas atraviesa toda la matriz, capturar el foco exige aumentar la resolución global y fragmentar zonas que no tienen patrón. El K-means de ST-CC reagrupa bloques ya formados, pero no divide ninguno por dentro, y elegir $K_R$ y $K_T$ con una curva o con Cross-Associations sigue dando una rejilla global.
\end{guion}

\begin{partes}
(a) Datos: un foco local en un bloque & El panel de la izquierda muestra los datos reales. El cuadro naranja es una zona pequeña con mucha actividad. & 2 regiones activas en 2 horarios, dentro de un bloque de 3 regiones por 4 horarios. \\
(b) Tablero de ST-CC: el foco queda promediado & El tablero pinta el bloque de un solo color y el foco desaparece. & Si 4 celdas valen 70 y 8 valen 10, el promedio es 30. \\
Resolución global & El número de grupos del tablero para toda la matriz. Subirlo parte todos los bloques. & Pasar de $3\times3$ a $4\times4$ parte también bloques sin patrón. \\
Fragmentar zonas sin patrón & Dividir en trozos pequeños partes de la matriz que eran uniformes. & Un bloque tranquilo de madrugada queda en cuatro pedazos inútiles. \\
K-means \se{ka-mins} & Técnica que agrupa objetos parecidos. ST-CC la usa para agrupar bloques ya formados por su promedio y dispersión. & Junta dos bloques que tienen actividad similar. \\
No subdivide un bloque por dentro & K-means junta bloques, nunca los divide. & Un bloque con foco sigue siendo un solo bloque. \\
Curva (elegir $K_R$ y $K_T$) & Gráfico de la función objetivo que se mira para elegir cuántos grupos usar. & Se busca el «codo» de la curva. \\
Cross-Associations \se{cros-asociéishons} & Método de Chakrabarti et al. (2004) que decide automáticamente cuántos grupos de filas y de columnas usar. & Elige $K_R$ y $K_T$, pero el resultado sigue siendo un tablero. \\
\end{partes}

\begin{preg}{¿Por qué no simplemente aumentar $K_R$ y $K_T$?}
Porque cada corte nuevo atraviesa toda la matriz y fragmenta también las zonas sin patrón. En AST-CC cada corte afecta solo al bloque donde se detectó heterogeneidad.
\end{preg}

% =====================================================================
\section{Planteamiento del problema}

\begin{guion}
La pregunta es cómo refinar localmente, y de forma estadísticamente justificada, la partición de ST-CC, manteniendo la interpretabilidad. Hay tres dificultades. No hay criterio para saber dónde refinar. Evaluar muchas divisiones sobre los mismos datos produce refinamientos espurios si no se controla la multiplicidad. Y la matriz suma varios días con una variabilidad mayor que la de Poisson, de modo que con millones de viajes casi todo resulta significativo.
\end{guion}

\begin{partes}
Refinar localmente & Dividir solo el bloque que lo necesita, sin tocar el resto. & Partir el bloque con el foco y dejar los otros ocho como están. \\
Estadísticamente justificada & Que cada división se decida con una prueba de hipótesis y no a ojo. & «Se divide porque la prueba dice que hay patrón». \\
Interpretabilidad & Que el resultado se pueda leer como grupos de regiones y de horarios. & «Estas 2 regiones se activan en estos 2 horarios». \\
¿Dónde refinar? & No hay un criterio que diga si un bloque tiene heterogeneidad real. & ¿Este bloque tiene un patrón o solo variación normal? \\
Heterogeneidad real & Que dentro del bloque no todo se comporta igual. & Regiones de oficina de día y regiones nocturnas dentro del mismo bloque. \\
¿Cuándo detenerse? & Si se divide sin control, siempre se encuentra algo y nunca se termina. & Hay que decidir cuándo ya no vale la pena dividir más. \\
Multiplicidad & Al hacer muchas pruebas sobre los mismos datos, crece la probabilidad de algún falso positivo. & Con 20 pruebas al 5\,\%, hay 64\,\% de probabilidad de al menos un falso rechazo. \\
Refinamientos espurios & Divisiones que parecen reales pero son pura casualidad. & Partir un bloque por un patrón que no existe. \\
¿Qué es ruido? & Variación aleatoria que no es un patrón real. & Un día llueve y baja el número de viajes en toda una zona. \\
Variabilidad mayor que la de Poisson \se{poasón} & Poisson supone que la variación entre días es igual al promedio. En los datos es mucho mayor (sobredispersión). & Un promedio de 100 viajes que varía como si fuera de 300. \\
Con millones de viajes casi todo resulta significativo & Con muchos datos, hasta diferencias minúsculas parecen significativas. & Una diferencia de 0.05\,\% sale «significativa» con $10^7$ viajes. \\
\end{partes}

\begin{preg}{¿Qué es un p-valor?}
Es la probabilidad, si no hubiera ningún patrón, de ver un resultado tan extremo como el observado. Un p-valor pequeño indica que es poco probable que sea casualidad.
\end{preg}

\begin{preg}{¿Cuál es la hipótesis nula en tu método?}
Que dentro del bloque todas las regiones tienen el mismo perfil horario y solo difieren en volumen. Rechazarla significa que hay un patrón espacio-temporal local.
\end{preg}

% =====================================================================
\section{Objetivos}

\begin{guion}
El objetivo general es proponer y validar una extensión de ST-CC basada en refinamiento jerárquico por bloques, guiado por pruebas de significancia derivadas de la I-divergencia. Los específicos son cuatro. OE1, formular el refinamiento como un árbol de bloques anidados. OE2, diseñar la regla de decisión y de parada. OE3, implementar y evaluar calibración y potencia con datos sintéticos y nulos. OE4, aplicarlo a los taxis de Manhattan y compararlo con ST-CC y rejillas globales.
\end{guion}

\begin{partes}
Objetivo general & Lo que la tesis quiere lograr en una frase: proponer y validar la extensión de ST-CC. & Es la «promesa» de toda la tesis. \\
Refinamiento jerárquico por bloques & Dividir bloques en bloques más pequeños, y estos a su vez en otros, formando niveles. & Un bloque se divide en dos y uno de esos se divide otra vez. \\
Pruebas de significancia derivadas de la I-divergencia & Pruebas de hipótesis construidas a partir de la misma medida de error que usa ST-CC. & El criterio para dividir es coherente con lo que optimiza el método. \\
OE1 \se{o e uno} & Formular el refinamiento como un árbol de bloques anidados, empezando por la partición base de BBAC$_I$. & Cada bloque base es la raíz de un pequeño árbol. \\
Árbol de bloques anidados & Estructura donde cada bloque puede tener «hijos». Los bloques finales (hojas) cubren toda la matriz. & Raíz $\rightarrow$ dos hijos $\rightarrow$ un hijo vuelve a dividirse. \\
Criterio de heterogeneidad & La medida que dice si un bloque merece dividirse. & Es $G^2_B$ («G cuadrado de B»). \\
OE2 \se{o e dos} & Diseñar la regla de decisión y la de parada. & Es la diapositiva 9 y la 10. \\
Prueba global por bloque & Una prueba sobre todo el bloque que dice si hay algo que refinar. & Se llama prueba ómnibus. \\
Mejor corte & Entre todas las formas posibles de dividir el bloque, la que más separa. & Entre 7 cortes de filas, el que más diferencia los dos grupos. \\
Calibración con la variabilidad entre días & Ajustar las pruebas usando cuánto varían realmente los datos de un día a otro. & Se remuestrean los días reales. \\
Control de multiplicidad & Repartir el nivel $\alpha$ para limitar los falsos refinamientos a lo largo de todo el árbol. & Cada bloque recibe una parte de $\alpha$. \\
Umbral de tamaño de efecto & Que el patrón, además de ser significativo, sea suficientemente grande. & Es el $\delta$ (delta) de la regla de decisión. \\
OE3 \se{o e tres} & Implementar el método y evaluar su calibración (que no invente patrones) y su potencia (que encuentre los reales). & Se prueba con datos sintéticos y con datos nulos. \\
Datos sintéticos & Datos fabricados donde se planta un patrón conocido. & Un foco que se sabe dónde está. \\
Datos nulos & Datos fabricados sin ningún patrón. & Sirven para ver si el método inventa cosas. \\
OE4 \se{o e cuatro} & Aplicar el método a los taxis de Manhattan y compararlo con ST-CC y con rejillas globales. & Se mide error de reconstrucción por parámetro, estabilidad e interpretabilidad. \\
Error de reconstrucción por parámetro & Qué tan bien reproduce la matriz el modelo, descontando cuántos parámetros usa. & Dos modelos con el mismo ajuste: gana el que usa menos parámetros. \\
Rejillas globales & Tableros más finos de ST-CC, que dividen toda la matriz. & Un tablero de $6\times6$ en lugar de $3\times3$. \\
\end{partes}

\begin{preg}{¿Cuál es el aporte principal?}
Un procedimiento que decide, bloque por bloque y con una prueba de hipótesis, si conviene subdividir. Controla el error a lo largo del árbol y entrega hojas que se leen como grupos de regiones y horarios.
\end{preg}

% =====================================================================
\section{Estado del arte: tres familias de trabajos desde 2020}

\begin{guion}
Revisé tres familias. La primera es co-agrupamiento espacio-temporal y bloques flexibles. Hay tableros globales o pruebas aplicadas como filtro posterior, y ninguno decide por bloque cuándo subdividir. La segunda es inferencia selectiva, que valida estructuras ya estimadas pero no construye un refinamiento local. La tercera es detección de \emph{hotspots}, que busca zonas de alta intensidad en el espacio y no interacciones dentro de bloques de una matriz región por tiempo. Parto de ST-CC para la partición base y de SigTree para el principio de construir con pruebas y detenerse cuando ninguna división es significativa.
\end{guion}

\begin{partes}
Co-agrupamiento y bloques flexibles & Trabajos que agrupan filas y columnas, o que relajan la rejilla fija. & Wu (2022), Neves (2022), Yen (2020), Sugahara (2023), Watanabe (2021). \\
Co-clusters coherentes \se{co-clusters} & Grupos de filas y columnas cuyos valores siguen la misma tendencia, aunque no sean iguales. & Dos estaciones cuya temperatura sube y baja igual. \\
\emph{Biclustering} \se{báiclasterin} & Buscar submatrices (algunas filas y algunas columnas) que se comportan de forma parecida, sin cubrir toda la matriz. & Neves (2022) lo usa en tráfico vial. \\
Modelos de bloques bayesianos y jerárquicos & Modelos probabilísticos que deciden el número de grupos o los organizan en niveles. & Yen (2020), Sugahara (2023). \\
Diferencia (familia 1) & Son tableros globales, o pruebas aplicadas como filtro posterior. Ninguno decide por bloque cuándo subdividir. & Es lo que hace único a AST-CC. \\
Inferencia selectiva & Pruebas que siguen siendo válidas después de haber agrupado los datos con esos mismos datos. & Gao (2024), Chen (2023), Watanabe (2021). \\
\emph{Count splitting} \se{caunt splíting} y \emph{data thinning} \se{deita zínin} & Técnicas que dividen los datos en dos partes: una para encontrar la estructura y otra para probarla. & Neufeld (2024). \\
Diferencia (familia 2) & Validan estructuras ya estimadas, pero ninguno construye un refinamiento jerárquico local. & AST-CC sí construye y prueba a la vez. \\
Hotspots \se{jótspots} & Zonas donde se concentra mucha actividad. & Una esquina con muchas paradas de taxi. \\
\emph{Space-time NetScan} \se{espéis-táim nétscan} & Método de Shiode (2022) que barre la red de calles para detectar focos en espacio y tiempo. & Detecta dónde y cuándo se concentran eventos. \\
Diferencia (familia 3) & Buscan zonas de alta intensidad, no interacciones dentro de bloques de una matriz regiones por tiempo. & Un foco no es lo mismo que un patrón interno de un bloque. \\
SigTree \se{sig-trí} & Método de Hu (2025) que construye un árbol con pruebas de hipótesis y se detiene cuando ninguna división es significativa. & De ahí viene la idea de parada de AST-CC. \\
Puntos de partida & ST-CC aporta la partición base; SigTree aporta el principio de construir con pruebas. & AST-CC une ambos. \\
\end{partes}

\begin{preg}{¿Qué es inferencia selectiva?}
Es hacer pruebas válidas cuando los grupos se eligieron con los mismos datos que luego se usan para probarlos. Si no se corrige, se detectan patrones falsos porque los grupos ya se escogieron para verse distintos.
\end{preg}

\begin{preg}{¿Por qué no usas \emph{count splitting} o \emph{data thinning}?}
Porque necesitan conocer la distribución de los conteos, y los conteos de viajes varían más de lo que permite Poisson. Separar la selección de la prueba dividiendo días queda como trabajo futuro.
\end{preg}

% =====================================================================
\section{Brecha de investigación}

\begin{guion}
La tabla compara los trabajos en cinco características. Las piezas existen por separado. ST-CC tiene datos espacio-temporales y filas con columnas, pero ni estructura local, ni criterio estadístico, ni control de error. SigTree tiene estructura local, criterio y control, pero no trabaja con datos espacio-temporales ni con filas y columnas a la vez. No encontré un trabajo que refine localmente los bloques de un co-agrupamiento espacio-temporal, decida cada división con una prueba de hipótesis y controle el error a lo largo de la jerarquía.
\end{guion}

\begin{partes}
Esp.-temp. \se{espacio-temporal} & Si el método trabaja con datos que tienen espacio y tiempo. & ST-CC sí; SigTree no. \\
Filas y col. \se{filas y columnas} & Si agrupa filas y columnas a la vez. & ST-CC sí; Gao y Shiode no. \\
Estructura local & Si permite refinar por zonas y no solo con cortes que cruzan todo. & SigTree, Neves, Gao y Shiode sí. \\
Criterio estad. \se{criterio estadístico} & Si decide o valida con una prueba de hipótesis. & ST-CC no; SigTree sí. \\
Control de error & Si controla los falsos positivos al hacer muchas pruebas. & Neves no; los demás con criterio estadístico sí. \\
$\checkmark$ \se{cumple} y $\times$ \se{no cumple} & Marcas que indican si el trabajo tiene esa característica. & Solo la fila de la propuesta tiene cinco $\checkmark$. \\
Brecha & El hueco que la tesis llena: refinar localmente, decidir con prueba de hipótesis y controlar el error. & «No se encontró» un trabajo que reúna las tres. \\
\end{partes}

\begin{preg}{¿Cómo sabes que no existe un trabajo así?}
Es el resultado de la revisión que hice, no una afirmación absoluta. Por eso la diapositiva dice «no se encontró». Las piezas existen por separado y la tesis las reúne.
\end{preg}

\begin{preg}{¿Neves (2022) ya usa una prueba de significancia?}
Sí, pero como filtro posterior sobre candidatos ya generados, y sin control de error en la jerarquía.
\end{preg}

% =====================================================================
\section{AST-CC: refinar mientras haya heterogeneidad detectable y relevante}

\begin{guion}
Mi propuesta se llama AST-CC. Parte del tablero de ST-CC, panel (a), y refina solo los bloques que lo necesitan, panel (b). Con conteos, la I-divergencia que optimiza BBAC$_I$ coincide con la devianza de Poisson. Lo que ST-CC no modela es la interacción dentro de cada bloque, que mido con $G^2_B$. Cuando corto un bloque, ese $G^2_B$ se reparte exactamente entre lo que explica el corte y lo que queda dentro de los hijos, así que cada corte se puede decidir con una prueba de hipótesis local. El método tiene cuatro etapas: armar la matriz, obtener la partición base, refinar cada bloque con pruebas y resumir las hojas.
\end{guion}

\begin{partes}
Título: heterogeneidad detectable y relevante & Detectable: que una prueba estadística la reconozca. Relevante: que sea suficientemente grande como para importar. & Con muchos datos todo es detectable, pero no todo es relevante. \\
(a) Partición base (ST-CC) & Tablero de $3\times3$ bloques sobre 9 regiones por 12 intervalos. Cada bloque tiene un solo color. & El foco queda diluido en el bloque de arriba a la derecha. \\
(b) Teselación refinada (AST-CC) \se{teselación} & Cobertura de la matriz con rectángulos que no se superponen. Aquí solo el bloque con el foco se divide en tres rectángulos. & Los otros ocho bloques no cambian. \\
Foco & La zona pequeña de alta actividad (el cuadro naranja). & Se aísla en su propia hoja. \\
Hojas & Los bloques finales que ya no se dividen. & Los tres rectángulos del bloque refinado. \\
Idea central & Resumen del método dentro de la caja naranja. & Se explica frase por frase abajo. \\
Con conteos, la I-divergencia coincide con la devianza de Poisson & La medida de error de BBAC$_I$ equivale a una medida estándar de ajuste para conteos. & Gracias a eso se pueden usar pruebas de hipótesis. \\
El término que ST-CC no modela & ST-CC resume cada bloque con un valor y no mira lo que ocurre dentro. & Regiones de mañana y de noche dentro del mismo bloque. \\
Interacción dentro del bloque & Que el perfil horario dependa de la región. & Unas regiones suben en la mañana y otras en la noche. \\
$G^2_B$ \se{G cuadrado de B} & Número que mide cuánta interacción hay dentro del bloque $B$. Vale cerca de 0 si todo es uniforme y crece cuando hay patrón. & 34 en el bloque de ejemplo. \\
Cortar un bloque la reparte exactamente & Si divido el bloque, la heterogeneidad total se reparte entre lo que explica el corte y lo que queda en los hijos, sin perder ni duplicar nada. & $34.1\approx33.6+0.25+0.25$ en el bloque de ejemplo. \\
$G^2_{B,s}$ \se{G cuadrado de B, s} & Parte del patrón que explica el corte $s$. & 33.6 en el ejemplo. \\
$s$ \se{ese} & Un corte concreto del bloque. & Separar las regiones $\{1,2\}$ de las $\{3,4\}$. \\
$C$ \se{ce} & Un bloque hijo, es decir, cada uno de los trozos que quedan al cortar. & Un corte de filas deja dos hijos. \\
$\sum_C G^2_C$ \se{sumatoria sobre C de G cuadrado de C} & Lo que queda de patrón dentro de los hijos, sumado. & $0.25+0.25$ en el ejemplo. \\
La fórmula completa \se{G cuadrado de B es igual a G cuadrado de B, s, más la suma sobre C de G cuadrado de C} & Patrón total $=$ lo que explica el corte $+$ lo que queda en los hijos. & $34.1=33.6+0.5$ (aprox.). \\
Cuatro etapas & Armar la matriz, obtener la partición base, refinar y resumir. & Ver la siguiente tabla. \\
\end{partes}

\subsection*{Las cuatro etapas}

\begin{partes}
1. Matriz \se{O, y X sub uno hasta X sub D} & $O$ es la matriz de conteos y $X_1,\dots,X_D$ son las matrices de cada día. Sumadas dan $O$. & 12 días laborables dan 12 matrices. \\
2. Partición base \se{BBAC sub I con K sub R, K sub T} & Se ejecuta BBAC$_I$ con $K_R$ grupos de filas y $K_T$ de columnas. & Resultado: el tablero. \\
3. Refinamiento & Cada bloque pasa por pruebas de hipótesis para decidir si se divide. & Es la novedad. \\
4. Resumen & Cada hoja se describe con su tamaño de efecto (cuánto patrón tiene) y su actividad (cuánto se mueve respecto de lo normal). & «Estas 2 regiones se activan a esta hora». \\
\end{partes}

\begin{preg}{¿Qué es $G^2_B$?}
Es el número que mide cuánto patrón interno tiene el bloque. Compara lo observado con lo que habría si la región no influyera en el horario. Vale cerca de 0 si el bloque es uniforme y crece cuando hay un patrón local.
\end{preg}

\begin{preg}{¿Qué es $C$?}
Un bloque hijo, es decir, cada uno de los trozos en que queda dividido el bloque después del corte.
\end{preg}

\begin{preg}{¿Qué es $G^2_{B,s}$?}
La parte del patrón que explica el corte $s$. Cuanto mayor es, más justifica el corte.
\end{preg}

\begin{preg}{¿Qué parte del error de ST-CC ataca AST-CC?}
La interacción que queda dentro de los bloques. ST-CC la deja sin modelar y AST-CC va trasladándola a hojas explicadas, hasta que lo que queda ya no se distingue del ruido.
\end{preg}

% =====================================================================
\section{Regla de decisión para cada bloque}

\begin{guion}
Cada bloque entra al diagrama con su nivel $\alpha_B$. Primero la puerta ómnibus: si la prueba global no es significativa o el tamaño de efecto es menor que $\delta$, es una hoja homogénea. Si pasa, se prueba si un corte de filas o de columnas es significativo y se aplica el de menor p-valor. Si ninguno lo es, se prueba el corte conjunto $2\times2$ como último recurso. Si tampoco, queda como hoja heterogénea no localizada. Un corte solo se acepta si su ganancia supera $\delta$, y entonces se recurre en cada hijo.
\end{guion}

\begin{partes}
Bloque $B$ con nivel $\alpha_B$ \se{bloque B con nivel alfa sub B} & El bloque que se está examinando y la cantidad de error que se le permite. $\alpha_B$ es la parte de $\alpha$ (el error total permitido) que le toca a este bloque, según cuántos viajes tiene. & Un bloque con 4000 de 10\,000 viajes recibe el 40\,\% del $\alpha$. \\
Ómnibus \se{ómnibus} & Prueba global sobre todo el bloque. Dice si hay algo que refinar, pero no dónde. «Ómnibus» significa «para todo». & Pregunta: ¿hay interacción región-horario en este bloque? \\
Significativo & Que el resultado es poco probable que sea casualidad, es decir, que el p-valor es menor que el nivel $\alpha_B$. & $p=0.002$ con $\alpha_B=0.01$ es significativo. \\
p-valor \se{pe valor} & Probabilidad de ver un patrón tan fuerte como el observado si en realidad no hubiera ninguno. Cuanto más pequeño, más convincente. & $p=0.2$ es casualidad probable. $p=0.002$, muy poco probable. \\
$\tilde I_B$ \se{I tilde de B} & Tamaño de efecto del bloque. Mide cuánto patrón hay, en una escala que no depende del número de viajes, y ya descontado el ruido (por eso lleva tilde). Se mide en nats por observación. & 0.03 es un patrón apreciable y 0.0005 es casi nada. \\
Nats \se{nats} & Unidad de información. Es la unidad con la que se mide la información mutua. & Como los metros para una longitud. \\
$\delta$ \se{delta} & Umbral mínimo de tamaño de efecto. Si $\tilde I_B$ no lo alcanza, el patrón se considera irrelevante aunque sea significativo. & Con $\delta=0.005$, un bloque con $\tilde I_B=0.002$ no se divide. \\
$\tilde I_B\ge\delta$ \se{I tilde de B mayor o igual que delta} & Pregunta de si el patrón es suficientemente grande. & $0.03\ge0.005$: sí. \\
Hoja homogénea & Bloque que ya no tiene nada que refinar. Es uniforme. & La ómnibus no rechaza. \\
Corte de filas o de columnas & Dividir el bloque en dos: separando regiones (filas) o separando horarios (columnas). & Separar regiones de mañana de regiones de noche. \\
Corte 1D \se{corte uno-D, o corte unidimensional} & Corte en una sola dimensión: solo filas o solo columnas. & Es lo contrario del corte conjunto. \\
Corte de menor p-valor \se{corte de menor pe valor} & Si hay varios cortes significativos, se elige el que tiene el p-valor más pequeño, es decir, el más convincente. & Entre un corte con $p=0.004$ y otro con $p=0.3$, se elige el primero. \\
Corte conjunto $2\times2$ \se{corte conjunto dos por dos} & Corta filas y columnas a la vez y produce cuatro hijos. Se usa solo si ningún corte de una dimensión es significativo. & Un patrón repartido entre varias regiones y varios horarios. \\
Último recurso & El corte conjunto se prueba solo cuando fallan los otros, porque añade más parámetros. & Se intenta al final. \\
Hoja heterogénea no localizada & Hay patrón en el bloque, pero ningún corte lo captura. Sugiere una interacción difusa. & La ómnibus rechaza y ningún corte pasa. \\
$\tilde g_{B,s}$ \se{g tilde de B, s} & Ganancia del corte $s$: cuánta heterogeneidad explica, también corregida por ruido y medida igual que $\tilde I_B$. & 0.08 significa que el corte explica bastante. \\
$\tilde g_{B,s}\ge\delta$ \se{g tilde de B, s mayor o igual que delta} & Pregunta de si el corte elegido vale la pena. & Un corte significativo que explica poco no se acepta. \\
Recurrir en cada hijo & Repetir todo el procedimiento en cada trozo resultante, con su propio nivel. & Cada hijo vuelve a pasar por el diagrama. \\
Un corte por nivel & En cada nivel se aplica un solo corte, no dos a la vez. & Se divide por filas ahora y por columnas, si hace falta, en el siguiente nivel. \\
\end{partes}

\subsection*{Cómo seguir el diagrama paso a paso}

\begin{center}
\small
\begin{tabular}{@{}p{7.0cm}p{8.0cm}@{}}
\toprule
Si el bloque tiene… & Resultado \\
\midrule
$p_{\text{ómnibus}}=0.2$ & Hoja homogénea (no hay evidencia de patrón). \\
Ómnibus significativa, $\tilde I_B=0.03$, corte de filas con $p=0.004$ y de columnas con $p=0.3$ & Se corta por filas y, si la ganancia supera $\delta$, se recurre en cada hijo. \\
Ómnibus significativa pero ningún corte (ni de una dimensión ni conjunto) es significativo & Hoja heterogénea no localizada. \\
Corte significativo pero ganancia menor que $\delta$ & Se rechaza el corte y queda como hoja heterogénea no localizada. \\
\bottomrule
\end{tabular}
\end{center}

\begin{preg}{¿Qué es la prueba ómnibus?}
Es una prueba global sobre todo el bloque. Dice si existe algo que refinar, pero no dónde. Las pruebas de cada corte localizan la estructura.
\end{preg}

\begin{preg}{¿Qué es el tamaño de efecto $\tilde I_B$ y por qué hace falta?}
Es cuánto patrón hay, sin depender de cuántos viajes tiene el bloque. Hace falta porque con muchos viajes casi todo es significativo, y el tamaño de efecto separa lo importante de lo trivial.
\end{preg}

\begin{preg}{¿Qué es $\delta$?}
El umbral mínimo de tamaño de efecto. Es la respuesta a «¿qué tan grande debe ser un patrón para que valga la pena refinar?». No se fija de antemano: se prueba con varios valores en el análisis de sensibilidad.
\end{preg}

\begin{preg}{¿Qué significa 1D?}
Una dimensión. Un corte 1D divide solo por filas (regiones) o solo por columnas (horarios). El corte $2\times2$ usa las dos a la vez.
\end{preg}

\begin{preg}{¿Por qué un solo corte por nivel?}
Cortar filas y columnas a la vez pasaría de uno a cuatro bloques cuando quizá bastan dos. Además, cortar una dimensión por nivel da una lectura directa: si se divide por filas, el patrón es espacial; si es por columnas, es temporal.
\end{preg}

\begin{preg}{¿Qué diferencia hay entre una hoja homogénea y una heterogénea no localizada?}
En la homogénea no hay patrón. En la heterogénea no localizada sí hay patrón, pero ningún corte simple lo captura. Se reportan por separado para no presentar como uniforme un bloque que no lo es.
\end{preg}

% =====================================================================
\section{Calibración, multiplicidad y parada}

\begin{guion}
Para calibrar las pruebas uso un \emph{bootstrap} de días. La varianza entre días es 3 veces la media en días laborables y 17 veces en feriados, así que remuestreo los días completos, lo que conserva la correlación entre celdas. Como el corte se elige porque separa bien, repito la búsqueda en cada réplica nula con un ajuste \emph{min-P}. Sin eso el error de tipo I llegó a 0.98. El nivel se reparte por masa a lo largo del árbol, lo que acota el error por familia por $\alpha$. Y se para por significancia y tamaño de efecto, que no crece con el número de viajes. El gráfico muestra que el \emph{bootstrap} de días es el que queda más cerca del nivel nominal de 0.05.
\end{guion}

\subsection*{Columna izquierda: los cuatro puntos}

\begin{partes}
Calibración & Ajustar las pruebas para que su p-valor sea confiable con datos reales. & Una prueba calibrada rechaza el 5\,\% de las veces cuando se le pide 0.05. \\
\emph{Bootstrap} de días \se{bútstrap de días} & Técnica de remuestreo. Se arman datos de prueba sorteando días completos de los datos reales, con reemplazo. & Con 12 días, se sortean 12 días (algunos pueden repetirse). \\
Varianza entre días & Cuánto varía una celda de un día a otro. & 3 veces su promedio en laborables y 17 veces en feriados. \\
Conserva la correlación entre celdas & Al usar días completos, las celdas de un mismo día se mueven juntas, como en la realidad. & Un día de lluvia baja a toda la zona al mismo tiempo. \\
Ajuste \emph{min-P} \se{mín pe} & Corrección por haber elegido el mejor corte entre muchos. Para cada dato de prueba se repite la búsqueda completa del mejor corte. & Es como comparar el mejor de 7 candidatos contra el mejor de 7 candidatos al azar. \\
Westfall (1993) & Autor del ajuste \emph{min-P} que se cita. & Referencia de la diapositiva. \\
Error de tipo I \se{error de tipo uno} & Declarar que hay patrón cuando no lo hay (falso positivo). & «Dividí el bloque, pero no había nada». \\
«El error de tipo I llegó a 0.98» & Sin el ajuste, casi todos los bloques sin patrón se dividían por error. & 98 de cada 100 falsos refinamientos. \\
Nivel por masa \se{nivel por masa} & El nivel de significancia se reparte según la cantidad de viajes de cada bloque. Los bloques con más datos reciben más. & Un hijo con 3/4 de los viajes recibe 3/4 del nivel del padre. \\
$\alpha_C=\alpha_B\,N_C/N_B$ \se{alfa sub C es igual a alfa sub B por N sub C sobre N sub B} & $\alpha_C$ es el nivel del hijo $C$, $\alpha_B$ el del padre, $N_C$ los viajes del hijo y $N_B$ los del padre. & Padre con 0.02 y un hijo con 3000 de 4000 viajes: $\alpha_C=0.015$. \\
Acota el error por familia por $\alpha$ \se{efe-doble u-e-erre} & La probabilidad de al menos un falso refinamiento en todo el árbol no supera $\alpha$. Se llama FWER. & Con $\alpha=0.05$, a lo sumo 5\,\% de riesgo de algún falso refinamiento. \\
Parada & Cuándo un bloque deja de dividirse. & Se detiene si no hay evidencia o el patrón es muy pequeño. \\
Significancia \emph{y} tamaño de efecto & Para dividir, el patrón debe ser significativo y también grande. & Las dos puertas de la diapositiva 9. \\
$\tilde I_B\ge\delta$ \se{I tilde de B mayor o igual que delta} & El tamaño de efecto no crece con el número de viajes, así que es una puerta justa para bloques grandes y pequeños. & $\tilde I_B=0.03$ pasa con $\delta=0.005$. \\
\end{partes}

\subsection*{Columna derecha: el gráfico}

El gráfico muestra el \textbf{error de tipo I al nivel 0.05} para cuatro maneras de calibrar la prueba. Cada método tiene dos barras: azul para días laborables y naranja para feriados. La línea punteada vertical marca el valor ideal (0.05).

\begin{partes}
Error de tipo I al nivel 0.05 & Porcentaje de veces que cada método declara patrón en datos que no tienen ninguno. Lo ideal es 0.05. & Una barra en 0.40 significa 40\,\% de falsas alarmas. \\
Datos nulos semirreales & Datos fabricados con los días reales pero sin patrón. & Sirven para medir falsas alarmas. \\
Media por bloque & Promedio del error de tipo I entre todos los bloques. & Un solo número por método y periodo. \\
Línea punteada \se{nominal cero cinco} & El valor al que se quiere llegar, 0.05. Se llama «nominal». & Cuanto más cerca de la línea, mejor. \\
Poisson & Supone que la variación es igual al promedio. Rechaza casi siempre porque los datos varían más. & 0.40 en laborables y 0.76 en feriados. \\
Cuasi-Poisson \se{cuasi-poasón} con $\hat\varphi$ constante \se{fi sombrero} & Poisson corregido con un factor único de sobredispersión para todo. & Es demasiado cauteloso en general: 0.03 y 0.01. \\
Gamma--Poisson \se{gama-poasón} con var $=a\mu^b$ \se{varianza igual a a por mu a la b} & Supone que la variación crece con el volumen siguiendo esa fórmula. & Da 0.02 y 0.06. \\
\emph{Bootstrap} de días & El método elegido. Es el que queda más cerca del 0.05. & 0.04 en laborables y 0.09 en feriados. \\
Laborable y feriado & Los dos periodos de los datos: días laborables y días feriados. & Azul y naranja. \\
Nota al pie: «algo liberal» & Con pocos días (feriado), el método rechaza un poco más de lo debido en bloques nocturnos. Es una limitación conocida. & «Liberal» significa que rechaza más de lo que debería. \\
\end{partes}

\begin{preg}{¿Qué es la sobredispersión?}
Que los datos varían de un día a otro mucho más de lo que permite el modelo de Poisson. Aquí la varianza es 3 veces el promedio en laborables y 17 veces en feriados.
\end{preg}

\begin{preg}{¿Qué es el error de tipo I y por qué importa?}
Es el falso positivo, es decir, dividir un bloque donde no hay patrón. Importa porque un método que inventa patrones no sirve. El objetivo es que ocurra a lo más con la frecuencia que indica el nivel ($\alpha$).
\end{preg}

\begin{preg}{¿Por qué no usas simplemente Poisson?}
Porque rechaza demasiado: en datos sin patrón declaró patrón en el 40\,\% de los bloques laborables y el 76\,\% de los feriados, cuando debería ser el 5\,\%.
\end{preg}

\begin{preg}{¿Por qué el \emph{bootstrap} es algo liberal en feriados?}
Porque solo hay 9 días y los bloques nocturnos son muy pequeños, de modo que la variación se estima con poca precisión. Es una limitación conocida y está declarada en la tesis.
\end{preg}

\begin{preg}{¿Qué es FWER?}
La probabilidad de cometer al menos un falso refinamiento en todo el árbol. Se controla repartiendo $\alpha$ entre los bloques.
\end{preg}

% =====================================================================
\section{Resultados esperados y forma de evaluación}

\begin{guion}
Evaluaré tres cosas. Primero, calibración: con datos nulos Poisson y semirreales, medir el error de tipo I y la FWER al nivel $0.05$, con la partición base estimada de los mismos datos. Segundo, potencia: con datos sintéticos con patrones locales plantados, medir tasa de detección, recuperación del patrón y falsos positivos. Tercero, datos reales de taxis de Manhattan en noviembre de 2017, comparados con ST-CC y rejillas globales con el mismo número de parámetros. Espero pruebas calibradas, patrones locales que ST-CC y las rejillas más finas no recuperan, menos parámetros para el mismo ajuste y hojas que se lean como grupos de regiones y horarios.
\end{guion}

\begin{partes}
OE3 Calibración \se{o e tres, calibración} & Comprobar que el método no inventa patrones. & Datos sin patrón: ¿cuántas divisiones hace? \\
Datos nulos Poisson y semirreales & Poisson: datos simulados sin patrón con la variación de Poisson. Semirreales: sin patrón pero con la variación de los días reales. & Dos niveles de dificultad para el método. \\
Partición base estimada de los mismos datos & El tablero inicial se calcula con los mismos datos que luego se prueban, que es la situación real. & Es la prueba más exigente. \\
Error de tipo I y FWER \se{error de tipo uno y efe-doble u-e-erre} & Falsa alarma por bloque y falsa alarma en todo el árbol, ambos al nivel 0.05. & Meta: no pasar de 5\,\%. \\
OE3 Potencia \se{o e tres, potencia} & Comprobar que el método sí encuentra los patrones que existen. & Se planta un foco y se mira si lo detecta. \\
Datos sintéticos con patrones locales plantados & Datos fabricados donde se sabe dónde está el patrón. & Foco de 2 regiones por 2 horarios. \\
Distinta intensidad & Se prueban patrones débiles y fuertes. & Un foco débil y otro fuerte. \\
Tasa de detección & Porcentaje de patrones plantados que el método encuentra. & 90 de 100 focos detectados. \\
Recuperación del patrón (ARI) \se{a erre i} & Qué tanto coincide la partición encontrada con la plantada. Vale 1 si son idénticas y cerca de 0 si coinciden como al azar. & ARI de 0.9 significa casi la misma partición. \\
Falsos positivos & Divisiones hechas donde no había nada plantado. & Cortes en bloques sin foco. \\
OE4 Datos reales & Aplicar el método a los taxis de Manhattan, noviembre de 2017, en un periodo laborable y otro feriado. & Los mismos datos de Liu et al. \\
Frente a ST-CC y rejillas globales con el mismo número de parámetros & Comparación justa: se compara con tableros que usan igual cantidad de parámetros. & Si AST-CC usa 400 parámetros, se compara con una rejilla que también usa 400. \\
Error de reconstrucción por parámetro (BIC) \se{bic} & Qué tan bien reproduce los datos el modelo, penalizando por cuántos parámetros usa. & Menor es mejor. \\
Fracción del término (iii) explicada \se{término tres} & Qué parte de la interacción interna que ST-CC deja sin modelar es explicada por el árbol. & Si es 60\,\%, el árbol explica 60\,\% de lo que ST-CC dejó. \\
Estabilidad por frecuencia de soporte & Se repite el método varias veces con distinto arranque y se cuenta cuántas veces reaparece cada hoja. & Una hoja que aparece en 9 de 10 corridas es estable. \\
Interpretabilidad (Penn Station) \se{pen estéishon} & Que las hojas tengan sentido urbano y se puedan leer directamente. & La zona de Penn Station, en Manhattan, es un ejemplo conocido. \\
Resultados esperados & Pruebas calibradas; patrones locales que ST-CC no recupera; menos parámetros para el mismo ajuste; hojas legibles. & Son las cuatro hipótesis de trabajo. \\
\end{partes}

\begin{preg}{¿Cómo evaluarás con datos reales si no conoces la verdad?}
Con tres lecturas: cuánto del término (iii) de ST-CC explica el árbol, si las hojas se repiten entre corridas con distinto arranque (frecuencia de soporte) y si tienen sentido urbano, por ejemplo la zona de Penn Station. Las pruebas con verdad conocida son las de datos sintéticos y nulos.
\end{preg}

\begin{preg}{¿Por qué comparar con el mismo número de parámetros?}
Porque cualquier modelo más flexible ajusta mejor. La pregunta justa es si refinar solo donde hay patrón logra el mismo ajuste con menos parámetros que dividir toda la matriz.
\end{preg}

\begin{preg}{¿Qué es el ARI?}
Una medida de qué tan parecidas son dos particiones, corregida por el azar. Vale 1 si son iguales. Aquí se usa para comparar las hojas encontradas con el patrón plantado.
\end{preg}

\begin{preg}{¿Qué es el BIC?}
Una medida de ajuste que penaliza el número de parámetros. Un modelo con más parámetros ajusta mejor, pero el BIC lo castiga. Sirve para comparar el árbol con tableros más finos en igualdad de condiciones.
\end{preg}

% =====================================================================
\section{Preguntas transversales difíciles}

\begin{preg}{¿El árbol que obtienes es el óptimo?}
No. Se construye de forma voraz, de arriba abajo, decidiendo cada corte con información local. Las decisiones son coherentes con el objetivo global, pero el árbol final no es necesariamente el mejor de todos los posibles.
\end{preg}

\begin{preg}{¿Qué pasa si BBAC converge a un mal punto de partida?}
La partición base puede variar entre ejecuciones. Por eso se corre con muchas inicializaciones y cada hoja se acompaña de su frecuencia de soporte. Una hoja que aparece en menos de la mitad de las ejecuciones no se interpreta como un hallazgo.
\end{preg}

\begin{preg}{¿Las pruebas son válidas si el tablero inicial se calculó con los mismos datos?}
Es el punto más delicado. La corrección con \emph{min-P} cubre la elección del corte dentro de un bloque, pero no la elección de los bloques base. Por eso el experimento de calibración con datos nulos usa la partición base estimada de los mismos datos y mide cuántos refinamientos falsos produce el método completo.
\end{preg}

\begin{preg}{¿Qué limitaciones reconoces?}
El \emph{bootstrap} es algo liberal en bloques nocturnos del feriado, donde hay pocos días. La garantía sobre el error por familia es aproximada, porque el árbol se construye con los mismos datos. Y la búsqueda del mejor corte es aproximada en bloques con más de 14 filas.
\end{preg}

\begin{preg}{¿Por qué significancia y tamaño de efecto a la vez?}
Porque la significancia crece con el número de viajes y el tamaño de efecto no. Con millones de viajes casi todo es significativo. El umbral $\delta$ exige además que el patrón sea suficientemente grande.
\end{preg}

\end{document}
